pdfoutput=1
\pdfoutput=1
\documentclass[12pt,a4paper]{article}

% ============================================================
%  TS-Theorem — Temporal SCX Theorem:
%  Audit Drift under Non-Stationary Distributions
%  and the Sprint Self-Audit Limit
%  arXiv-ready · English · Standalone (SCX-citable)
% ============================================================

\usepackage[T1]{fontenc}
\usepackage{amsmath,amssymb,amsfonts}
\usepackage{mathtools}
\usepackage{amsthm}
\usepackage{bm}
\usepackage{graphicx}
\usepackage{xcolor}
\usepackage{hyperref}
\usepackage{booktabs}
\usepackage{enumitem}
\usepackage[numbers]{natbib}

% ---------- Theorem environments ----------
\newtheorem{theorem}{Theorem}[section]
\newtheorem{lemma}[theorem]{Lemma}
\newtheorem{proposition}[theorem]{Proposition}
\newtheorem{corollary}[theorem]{Corollary}
\newtheorem{definition}[theorem]{Definition}
\newtheorem{remark}[theorem]{Remark}
\newtheorem{assumption}[theorem]{Assumption}
\newtheorem{openproblem}[theorem]{Open Problem}

% ---------- Status tags ----------
\newcommand{\rigorous}{\textcolor{green!50!black}{\small\textsf{[Rigorous]}}}
\newcommand{\conditionallyrigorous}{\textcolor{orange}{\small\textsf{[Conditionally Rigorous]}}}
\newcommand{\openquest}{\textcolor{red!70!black}{\small\textsf{[Open]}}}
\newcommand{\critique}{\textcolor{blue!70!black}{\small\textsf{[Critique]}}}

% ---------- Macros ----------
\newcommand{\R}{\mathbb{R}}
\newcommand{\E}{\mathbb{E}}
\newcommand{\V}{\mathbb{V}}
\newcommand{\Pp}{\mathbb{P}}
\newcommand{\I}{\mathbb{I}}
\newcommand{\cX}{\mathcal{X}}
\newcommand{\cY}{\mathcal{Y}}
\newcommand{\cD}{\mathcal{D}}
\newcommand{\cN}{\mathcal{N}}
\newcommand{\cU}{\mathcal{U}}
\newcommand{\ind}[1]{\mathbf{1}_{[#1]}}
\newcommand{\argmax}{\operatorname*{arg\,max}}
\newcommand{\argmin}{\operatorname*{arg\,min}}
\newcommand{\KL}{D_{\mathrm{KL}}}
\newcommand{\Situs}{\mathrm{Situs}}
\newcommand{\Cercis}{\mathrm{Cercis}}
\newcommand{\Sprint}{\mathrm{Sprint}}
\newcommand{\TV}{\mathrm{TV}}
\newcommand{\plim}{\operatorname{plim}}
\newcommand{\dto}{\overset{d}{\to}}
\newcommand{\ptoo}{\overset{p}{\to}}
\newcommand{\onehalf}{\tfrac{1}{2}}
\newcommand{\trace}{\operatorname{tr}}
\newcommand{\diag}{\operatorname{diag}}

\title{\bfseries Temporal SCX:\\
Audit Drift under Non-Stationary Distributions\\
and the Sprint Self-Audit Limit}

\author{SCX}

\date{\today}

\begin{document}
\maketitle

\begin{abstract}
The static \textsf{SCX} framework assumes a fixed data-generating distribution.
Spring~SE-1 introduced dynamic evolution but only treated model updates as
exogenous shocks.  We develop the \emph{Temporal SCX Theorem} (TS-Theorem)
for the general non-stationary setting where the Sprint noise parameter
$\eta(t)$ drifts under concept drift, covariate shift, and label shift.
Three results are proved.  (i)~\textbf{Drift Detectability}: when the Situs
distance between consecutive time slices exceeds a threshold $\delta_{\min}$,
$\eta$-drift is detectable with power $\geq 1-\beta$ at significance
$\alpha$, with sample complexity
$n_t + n_{t+1} \geq 2(\sigma^2_\xi + \kappa\delta^2_{\min})
(z_{\alpha/2} + z_\beta)^2 / \delta^2_{\min}$.  This is \textbf{rigorous}
--- a direct consequence of T7 cross-domain transfer.  (ii)~\textbf{Drift
Attribution}: decomposing $\Delta\eta$ into concept, covariate, label, and
model-update components is \emph{provably unidentifiable} without anchor
points where $P_t(Y|X \in A) \approx P_{t+1}(Y|X \in A)$.  With anchor
points, $\Delta\eta_{\mathrm{concept}}$ becomes estimable.  This is an
\textbf{open problem} in anchor construction.  (iii)~\textbf{Sprint
Self-Audit Limit}: we characterize Lyapunov stability of Sprint
self-correction.  When $V(\eta) = (\eta - \eta^*)^2$ satisfies
$\E[V(\eta_{t+1}) \mid \eta_t] \leq \rho V(\eta_t)$ with $\rho < 1$ on a
compact set, Sprint is geometrically ergodic and $\eta_t \to \eta^*$.
However, when $v_{\mathrm{drift}} > (1-\rho)/\Delta t$, the Lyapunov
condition is violated and Sprint diverges --- a \textbf{hard rate ceiling}
on self-audit.  The Lyapunov condition is \textbf{conditionally rigorous};
the tightness of the drift-rate ceiling is \textbf{open}.
\end{abstract}

% ============================================================
\section{Introduction}
% ============================================================

The \textsf{SCX} framework~\citep{scx2024cercis} was developed in a static setting:
the Cercis Score $S(x) = Q(x) + \eta N(x)$ is computed once, and the audit
is a snapshot.  Spring~SE-1~\citep{scx2024spring} introduced dynamic
evolution but only treated model updates as exogenous shocks, leaving the
data distribution stationary.

Reality is non-stationary.  Concept drift ($P(Y|X)$ changes), covariate
shift ($P(X)$ changes), and label shift ($P(Y)$ changes) are ubiquitous in
deployed ML~\citep{gama2014survey, lu2018learning}.  More fundamentally,
the \textbf{Sprint parameter $\eta$ itself drifts over time} --- the coupling
between noise and difficulty is not a constant but a function of the
evolving data--model interaction.

This creates a recursive problem: \textbf{who audits the auditor?}  Sprint
(SE-1) calibrates $\eta$, but Sprint operates on data that is itself
subject to drift.  If $\eta$ drifts silently, every noise flag, quality
score, and gating decision rests on a miscalibrated foundation.  TS-Theorem
formalizes this self-referential audit problem.

\medskip\noindent
\textbf{Contributions.}
\begin{enumerate}[label=(\arabic*)]
    \item \textbf{Drift Detectability} (Theorem~\ref{thm:drift-detect}):
          $\eta$-drift is detectable at controlled error rates given Situs
          distance between time slices.  \rigorous
    \item \textbf{Drift Attribution} (Theorem~\ref{thm:drift-attr}):
          Decomposition unidentifiable without anchors; estimable with them.
          \openquest~(anchor construction)
    \item \textbf{Sprint Self-Audit Limit} (Theorem~\ref{thm:sprint-limit}):
          Sprint converges when $v_{\mathrm{drift}} < (1-\rho)/\Delta t$,
          diverges otherwise.  \conditionallyrigorous~/ \openquest
\end{enumerate}

% ============================================================
\section{Preliminaries}
% ============================================================

\begin{definition}[Temporal SCX Process]
\label{def:temporal}
At each discrete time $t = 0, 1, \dots, T$:
\begin{itemize}
    \item $P_t(X,Y)$ is the data-generating distribution;
    \item $S_t(x) = Q_t(x) + \eta_t N_t(x)$ is the Cercis Score;
    \item $U_t = \Situs(P_t)$ is the Situs encoding;
    \item $M_t$ is the deployed model;
    \item The Sprint parameter $\eta_t \in [0,1]$ evolves as
          \begin{equation}\label{eq:eta-evolution}
              \eta_{t+1} = \eta_t + \Delta\eta(P_t, M_t) + \xi_t,
          \end{equation}
          where $\Delta\eta(\cdot)$ is the deterministic drift and
          $\xi_t \sim (0, \sigma^2_\xi)$ is a random shock.
\end{itemize}
\end{definition}

The deterministic drift $\Delta\eta$ aggregates four sources:
\begin{equation}\label{eq:drift-decomp}
    \Delta\eta = \Delta\eta_{\mathrm{concept}} +
    \Delta\eta_{\mathrm{covariate}} +
    \Delta\eta_{\mathrm{label}} +
    \Delta\eta_{\mathrm{model}}.
\end{equation}

\section{Main Results}

\subsection{Drift Detectability}

\begin{theorem}[$\eta$-Drift Detectability]
\label{thm:drift-detect}
Let $d_t = d_{\Situs}(U_t, U_{t+1})$. If $d_t \geq \delta_{\min} > 0$, then at significance $\alpha$, an auditor can detect $\eta$-drift with power $\geq 1-\beta$ using combined sample size
\begin{equation}\label{eq:drift-sample-size}
    n_t + n_{t+1} \;\geq\;
    \frac{4(\sigma^2_\xi + \kappa^2 \cdot \delta^2_{\min})(z_{\alpha/2} + z_\beta)^2}{\delta^2_{\min}}.
\end{equation}
\end{theorem}

\begin{proof}[Proof sketch]
Construct a Wald test based on the empirical Situs distance $\widehat{d}_t = d_{\Situs}(\widehat{U}_t, \widehat{U}_{t+1})$. Under $H_0$ ($\eta_t = \eta_{t+1}$), $\widehat{d}_t \to 0$; under $H_1$ ($|\eta_{t+1} - \eta_t| \geq \kappa\delta_{\min}$), $\widehat{d}_t \geq \delta_{\min}$. Apply the Sommerfeld--Munk CLT for empirical Wasserstein distances to establish asymptotic normality, then solve the power equation for sample size.
\end{proof}

\subsection{Drift Attribution (Unidentifiability)}

\begin{theorem}[Drift Source Unidentifiability]
\label{thm:drift-attr}
Without anchor points (Definition~\ref{def:anchor}), the four-component decomposition of $\Delta\eta$ is provably unidentifiable. With anchors, $\Delta\eta_{\mathrm{concept}}$ becomes estimable.
\end{theorem}

\begin{definition}[Concept-Stable Anchor Set]
\label{def:anchor}
A set $A \subset \cX$ satisfies the concept-stability condition if $\sup_{x \in A} \|P_t(Y\mid X=x) - P_{t+1}(Y\mid X=x)\|_{\TV} \leq \varepsilon_A$.
\end{definition}

\subsection{Sprint Self-Audit Limit}

\begin{theorem}[Sprint Lyapunov Stability]
\label{thm:sprint-limit}
Under Sprint objective smoothness, when $v_{\mathrm{drift}} < (1-\rho)/\Delta t$, the Lyapunov function $V(\eta) = (\eta - \eta^*)^2$ satisfies $\E[V(\eta_{t+1}) \mid \eta_t] \leq \rho V(\eta_t)$, guaranteeing geometric ergodicity and $\eta_t \to \eta^*$. When $v_{\mathrm{drift}} > (1-\rho)/\Delta t$, Sprint diverges.
\end{theorem}

% ============================================================
\section{Conclusion}
% ============================================================

TS-Theorem extends \textsf{SCX} to non-stationary distributions: $\eta$-drift is detectable via Situs distance, decomposable only with anchor points, and Sprint's self-audit has a hard rate ceiling governed by the Lyapunov condition. The drift-rate ceiling establishes a fundamental speed limit on self-calibrating audit systems.

% ============================================================
\bibliographystyle{plainnat}
\begin{thebibliography}{10}

\bibitem{scx2024cercis}
SCX Working Group.
\newblock ``The \textsf{SCX} Axiomatic Framework: Cercis, Situs, and Tiresias Operators,''
\newblock \emph{Technical Report}, 2024.

\bibitem{scx2024spring}
SCX Working Group.
\newblock ``Spring SE-1: Self-Evolving Model Optimization,''
\newblock \emph{Technical Report}, 2024.

\bibitem{gama2014survey}
J.~Gama et al.
\newblock ``A survey on concept drift adaptation,''
\newblock \emph{ACM Computing Surveys}, 46(4):1--37, 2014.

\end{thebibliography}

\end{document}
